% Dorian's notes and comments
% Personal scratch file --- not included in main.tex

\section*{Notes: Building a Simple Snake Game in Python}

This note explains a small Snake game written in Python using the
\texttt{pygame} library. The source file is \texttt{snake.py} and lives
outside the Overleaf repository (in the parent project folder).

\subsection*{What you need to run it}

Python 3 and the \texttt{pygame} library. Install pygame once with:

\begin{quote}
\texttt{pip install pygame}
\end{quote}

Then launch the game with \texttt{python snake.py}. Steer the snake with
the arrow keys or W/A/S/D. Eat the red squares to grow. Don't run into
the walls or yourself. Press \texttt{R} after a game over to restart.

\subsection*{The core idea: a grid of cells}

Instead of moving the snake pixel by pixel, the playing field is treated
as a grid of cells. Each cell is 20 pixels wide. The grid is 30 cells
across and 20 cells tall, which gives a window of $600 \times 400$
pixels. Everything in the game lives on this grid: the snake is a list
of cell coordinates, and the food is a single cell coordinate. Thinking
in grid units rather than pixels makes collision checks and movement
simple integer comparisons.

\subsection*{Representing the snake}

The snake is stored as a Python list of \texttt{(x, y)} tuples. The
first element is the head, the rest are the body, the last is the tail.
A direction is also stored as a tuple, for example \texttt{(1, 0)} means
``one cell to the right per step'' and \texttt{(0, -1)} means ``one cell
up per step''.

\subsection*{The game loop}

The whole game runs inside one big \texttt{while True} loop. Each pass
through the loop is one ``tick'' of the game and does four things in
order:

\begin{enumerate}
    \item \textbf{Handle input.} Read all keyboard events since the last
    tick. If the user pressed an arrow key, remember the new direction
    in a separate variable called \texttt{pending}. The active direction
    is only updated at the start of the next tick. This avoids a classic
    bug where pressing two keys quickly (for example right then down
    while moving right) could cause the snake to reverse into itself.
    A guard also rejects any key that would directly reverse the current
    direction.

    \item \textbf{Move the snake.} Compute a new head position by adding
    the direction vector to the current head. Then check for collisions
    (described below). If the move is legal, insert the new head at the
    front of the list. If the new head landed on the food, increase the
    score and place new food; otherwise, remove the tail. Adding a head
    and removing a tail makes the snake appear to slide forward one
    cell without rebuilding the whole list.

    \item \textbf{Draw.} Fill the screen black, then draw every snake
    segment as a green square and the food as a red square. The head is
    drawn slightly brighter than the body so you can see which way the
    snake is facing. The score is drawn in the top-left corner. If the
    game is over, a centred message is drawn over the field.

    \item \textbf{Pace the loop.} Call \texttt{clock.tick(FPS)} with
    \texttt{FPS = 10}, which forces the loop to run roughly ten times
    per second. Without this the loop would run as fast as the CPU
    allows and the snake would be uncontrollable.
\end{enumerate}

\subsection*{Collisions}

Two kinds of collisions end the game:

\begin{itemize}
    \item \textbf{Wall:} the new head coordinate falls outside the grid,
    that is its $x$ is not in $[0, 30)$ or its $y$ is not in $[0, 20)$.
    \item \textbf{Self:} the new head coordinate appears anywhere else
    in the snake's list of segments.
\end{itemize}

If either happens, a \texttt{game\_over} flag is set. The drawing code
keeps running so the final frame stays visible, and a key handler waits
for \texttt{R} to restart.

\subsection*{Placing food}

Food is placed by picking a random cell and checking it is not already
occupied by the snake. If it is, another random cell is picked. For a
small grid this rejection-sampling approach is fast enough and avoids
having to track every free cell separately.

\subsection*{Why this design}

Snake is a good first game because every interesting concept stays
simple:

\begin{itemize}
    \item State is just a list of tuples, an integer score, and a
    direction vector --- no classes or inheritance needed.
    \item Movement is one vector addition per tick.
    \item Collision is one bounds check and one membership test.
    \item The render step is a handful of coloured rectangles.
\end{itemize}

The whole game fits comfortably in around a hundred lines and is a
useful template for any other grid-based game (Tetris, Pac-Man, etc.):
keep a model of the world as plain data, advance it one tick at a time,
and redraw from scratch each frame.
